% formal/hko-local-maximality-conditions.tex — input by formal/main.tex
%
% Second-order sufficient conditions for local maximality of the systolic ratio.
%
% Sources:
%   - Generated data: experiments/hko-local-maximum/empirical/second-order/second-order-base.jsonl
%   - Gradient analysis: formal/hko-symmetry-gradient-structure.tex
%   - Capacity derivatives: formal/capacity-derivatives.tex [lem:cap-derivative]
%   - Volume derivatives: formal/capacity-derivatives.tex [lem:vol-derivative]

\subsection{Second-order analysis of flat directions}
\label{sec:second-order-analysis}

% [TODO: JÖRN - verify that this non-smooth second-order sufficiency argument
% is mathematically rigorous. The key claim (Proposition~\ref{prop:second-order-local-max})
% uses Danskin's theorem for the first-order part and a direct second-order argument
% for flat directions. The non-smooth analysis (orbit switching) may need additional
% justification from non-smooth optimization theory.]

\subsubsection{Setup}

Let $K = \{x \in \mathbb{R}^4 : a_i^\top x \leq 1,\; i = 1,\ldots,F\}$ be
a convex polytope with dual vertices $a_i \in \mathbb{R}^4$.
The systolic ratio is
\[
  \operatorname{sys}(K) = \frac{c_{\mathrm{EHZ}}(K)^2}{2\,\operatorname{vol}(K)}.
\]
% Citation: HK2017 Theorem 1.1 (formula_theorem) gives c_EHZ as the maximum of a QP
% for general polytopes. Note: HK2017 1.4 is a Remark, not a Theorem. The billiard
% characterization is from Artstein-Avidan \& Ostrover (2014), not HK2017.
% The KKT reformulation is developed in our library.
The EHZ capacity is the minimum action over all closed characteristics (Reeb orbits):
\[
  c_{\mathrm{EHZ}}(K) = \min_{i \in \mathcal{I}} c_i(K),
\]
where $\mathcal{I}$ indexes the finitely many combinatorial orbit types $(S,\sigma)$
and each $c_i(K)$ is the capacity computed from the KKT system for orbit type~$i$.

Each $c_i$ is a smooth function of the dual vertices $a_k$
(Lemma~\ref{lem:cap-derivative}), hence each per-orbit systolic ratio
$\operatorname{sys}_i(K) = c_i(K)^2 / (2\operatorname{vol}(K))$ is smooth.
The systolic ratio is
\begin{equation}\label{eq:sys-min-envelope}
  \operatorname{sys}(K) = \min_{i \in \mathcal{I}} \operatorname{sys}_i(K).
\end{equation}

\subsubsection{First-order necessary condition}

By Danskin's theorem, the directional derivative of $\operatorname{sys}$ at $K$
in direction $d \in \mathbb{R}^{4F}$ is
\begin{equation}\label{eq:danskin-directional}
  D_d^+ \operatorname{sys}(K) = \min_{i \in \mathcal{A}(K)} \nabla_{a} \operatorname{sys}_i(K) \cdot d,
\end{equation}
where $\mathcal{A}(K) = \{i \in \mathcal{I} : \operatorname{sys}_i(K) = \operatorname{sys}(K)\}$
is the active set of orbits achieving the minimum.

\begin{lemma}[First-order necessary condition]\label{lem:first-order-necessary}
  If $K$ is a local maximum of $\operatorname{sys}$, then
  $0 \in \operatorname{conv}\{\nabla_{a}\operatorname{sys}_i(K) : i \in \mathcal{A}(K)\}$.
\end{lemma}
\begin{proof}
  If $K$ is a local maximum, then $D_d^+ \operatorname{sys}(K) \leq 0$ for all $d$.
  By~\eqref{eq:danskin-directional}, this means
  $\min_{i \in \mathcal{A}} \nabla\operatorname{sys}_i \cdot d \leq 0$ for all $d$,
  which is equivalent to $0 \in \operatorname{conv}\{\nabla\operatorname{sys}_i\}$
  (standard separation argument).
\end{proof}

\subsubsection{Flat directions}

Let $G \in \mathbb{R}^{|\mathcal{A}| \times 4F}$ be the gradient matrix
with rows $\nabla_a \operatorname{sys}_i(K)^\top$ for $i \in \mathcal{A}(K)$.

\begin{definition}[Non-decreasing cone and flat directions]\label{def:flat-directions}
  The \emph{non-decreasing cone} is
  \[
    C = \{d \in \mathbb{R}^{4F} : \langle \nabla\operatorname{sys}_i, d \rangle \geq 0
    \text{ for all } i \in \mathcal{A}\},
  \]
  the intersection of $|\mathcal{A}|$ closed half-spaces through the origin.
  A direction $d$ is \emph{flat} if $D_d^+ \operatorname{sys}(K) = 0$,
  i.e., $d$ lies on the boundary of~$C$.
\end{definition}

\begin{lemma}[Non-decreasing cone collapses to kernel]\label{lem:cone-equals-kernel}
  Suppose $0 \in \operatorname{conv}\{\nabla\operatorname{sys}_i : i \in \mathcal{A}\}$
  with convex coefficients $\lambda_i \geq 0$, and let
  $\mathcal{A}^+ = \{i \in \mathcal{A} : \lambda_i > 0\}$.
  If $\operatorname{rank}(G_{\mathcal{A}^+}) = \operatorname{rank}(G)$,
  then $C = \ker(G)$.
  In particular, every flat direction satisfies
  $\langle \nabla\operatorname{sys}_i, d \rangle = 0$
  for \emph{all} $i \in \mathcal{A}$, and the set of flat directions
  is a linear subspace.
\end{lemma}
\begin{proof}
  $\ker(G) \subseteq C$ is immediate. For the reverse, let $d \in C$.
  Then $\langle \nabla\operatorname{sys}_i, d \rangle \geq 0$ for all $i$.
  Since $0 = \sum_i \lambda_i \nabla\operatorname{sys}_i$, we have
  $0 = \sum_i \lambda_i \langle \nabla\operatorname{sys}_i, d \rangle$.
  Each term is $\geq 0$ (product of $\lambda_i \geq 0$ and
  $\langle \nabla\operatorname{sys}_i, d \rangle \geq 0$), so each term
  equals~$0$. For $i \in \mathcal{A}^+$ ($\lambda_i > 0$), this gives
  $\langle \nabla\operatorname{sys}_i, d \rangle = 0$, hence
  $d \in \ker(G_{\mathcal{A}^+})$.
  Since $\operatorname{rank}(G_{\mathcal{A}^+}) = \operatorname{rank}(G)$,
  we have $\ker(G_{\mathcal{A}^+}) = \ker(G)$, so $d \in \ker(G)$.
\end{proof}

If $0 \in \operatorname{int}(\operatorname{conv}\{\nabla\operatorname{sys}_i\})$,
then $\ker(G) = \{0\}$ and $K$ is a strict first-order local maximum
(all directional derivatives are strictly negative).
Otherwise, $\ker(G)$ is nontrivial and second-order analysis is needed.

\subsubsection{Second-order sufficient condition}

\begin{proposition}[Second-order sufficient condition for local maximality]
\label{prop:second-order-local-max}
  Suppose $K$ satisfies the first-order necessary condition
  ($0 \in \operatorname{conv}\{\nabla\operatorname{sys}_i\}$),
  the rank condition of Lemma~\ref{lem:cone-equals-kernel}
  ($\operatorname{rank}(G_{\mathcal{A}^+}) = \operatorname{rank}(G)$,
  so that $C = \ker(G)$), and additionally:
  \begin{enumerate}
    \item For every non-flat direction $d \notin \ker(G)$: $D_d^+ \operatorname{sys}(K) < 0$
    (which follows from $C = \ker(G)$ and $Gd \neq 0$).
    \item For every flat direction $d \in \ker(G) \setminus \{0\}$:
    \begin{equation}\label{eq:second-order-condition}
      \limsup_{\varepsilon \to 0} \frac{\operatorname{sys}(K + \varepsilon d) - \operatorname{sys}(K)}{\varepsilon^2} < 0.
    \end{equation}
  \end{enumerate}
  Then $K$ is a strict local maximum of $\operatorname{sys}$.
\end{proposition}

% [TODO: JÖRN - the proof below is a sketch. The compactness argument
% on S^{4F-1} needs care because the convergence rate in condition (1) may
% degenerate as d approaches ker(G). A full proof would need uniform
% bounds, e.g., via the angle between d and ker(G).]
\begin{proof}[Proof sketch]
  Decompose any direction $d = d_\perp + d_\parallel$ where
  $d_\parallel \in \ker(G)$ and $d_\perp \perp \ker(G)$.

  For $d_\perp \neq 0$: by condition~(1),
  $D_{d_\perp}^+ \operatorname{sys} < 0$, so
  $\operatorname{sys}(K + \varepsilon d) \leq
  \operatorname{sys}(K) + \varepsilon D_{d_\perp}^+ \operatorname{sys} + O(\varepsilon^2)
  < \operatorname{sys}(K)$
  for small $\varepsilon > 0$.

  For $d_\perp = 0$ (pure flat): condition~(2) gives
  $\operatorname{sys}(K + \varepsilon d) < \operatorname{sys}(K)$
  for small $|\varepsilon|$.

  For mixed directions: the first-order decrease from $d_\perp$ dominates
  the second-order behavior along $d_\parallel$.

  By compactness of $S^{4F-1}$, the required $\varepsilon$ is bounded below
  uniformly, so $K$ is a strict local maximum.
\end{proof}

\begin{remark}\label{rem:non-smooth-curvature}
  The limit in~\eqref{eq:second-order-condition} is the curvature of the
  min-envelope~\eqref{eq:sys-min-envelope} along the flat direction~$d$.
  When the same orbit $i$ achieves the minimum for all $|\varepsilon|$ small,
  this equals $\frac{1}{2} d^\top \nabla^2 \operatorname{sys}_i(K) d$.
  At orbit-switching boundaries (where the minimizing orbit changes),
  the curvature may be discontinuous, but its sign remains well-defined.
\end{remark}

\subsubsection{Computational results for HKO2024}

For HKO2024 with $F = 10$, the gradient matrix $G$ has $|\mathcal{A}| = 150$
active orbits and lives in $\mathbb{R}^{150 \times 40}$
(dual-vertex parameterization, $4F = 40$ dimensions, no gauge).

\paragraph{First-order analysis.}
\begin{itemize}
  \item $\operatorname{rank}(G) = 25$, so $\dim\ker(G) = 15$.
  \item The LP test confirms $0 \in \operatorname{conv}\{\nabla\operatorname{sys}_i\}$
    with $26$ orbits having $\lambda_i > 0$ (Lemma~\ref{lem:first-order-necessary}).
  \item $\operatorname{rank}(G_{\mathcal{A}^+}) = 25 = \operatorname{rank}(G)$,
    so by Lemma~\ref{lem:cone-equals-kernel}, $C = \ker(G)$:
    the flat directions form a $15$-dimensional linear subspace.
\end{itemize}

\paragraph{Second-order analysis.}
\begin{itemize}
  \item For $15$ orthonormal basis vectors $d_0, \ldots, d_{14}$ of $\ker(G)$
    (from SVD), the curvature ratio
    $r(\varepsilon) = (\operatorname{sys}(K + \varepsilon d) + \operatorname{sys}(K - \varepsilon d) - 2\operatorname{sys}(K)) / \varepsilon^2$
    is negative for all $15$ directions, with median values ranging from
    $-0.31$ to $-0.02$ (second-order-curves.jsonl).
  \item Negative curvature on all basis vectors implies the diagonal of the
    generalized Hessian $H$ restricted to $\ker(G)$ is negative in the SVD
    basis, but does not prove negative definiteness (off-diagonal entries
    could produce positive eigenvalues).
  \item To test negative definiteness: $100$ random unit directions
    in $\ker(G)$ (uniform on $S^{14}$) all have negative curvature,
    with the least negative being $-0.06$ (second-order-random.jsonl).
\end{itemize}

Condition~(1) is verified exactly (it follows from the LP and rank
condition). Condition~(2) is checked numerically along $15$ basis
directions and $100$ random directions in $\ker(G)$ (all negative,
worst curvature $-0.06$). This does not constitute a proof of
condition~(2) for all $d \in \ker(G)$, but provides strong numerical
evidence that HKO2024 is a strict local maximum of $\operatorname{sys}$
among $F = 10$ polytopes in $\mathbb{R}^4$.
